\documentstyle[%


righttag%


,11pt%


,amssymb%


,russian%


,izvru%


]{amsart}





\newcommand{\tl}[3]{\medskip\noindent{\bf#1}\ #2 \dotfill #3 \par} 


\pagestyle{empty}


% \pagestyle{plain}


\begin{document}


% \setcounter{page}{80}


\par


\par


\par


\bigskip





\centerline{\Large СОДЕРЖАНИЕ}


\bigskip


\bigskip


%%%%%%%%%% Use \newline %%%%%%%%%%%











\tl{Аглямзянова Г.Н., Хайруллин Р.С.} 


{Задача Трикоми в классе функций, 


неограниченных\newline на характеристике}{3}


\tl{Аксентьев Л.А.} 


{Выпуклость поверхности конформного радиуса 


и оценки коэффициентов\newline отображающей функции}{8}


\tl{Апайчева Л.А.}


{Оптимальные квадратурные и кубатурные формулы для сингулярных\newline


интегралов с ядрами Гильберта}{16}


\tl{Кожевников В.В.}


{Об одной форме прямого вариационного метода


для однолистных\newline функций}{28}


\tl{Микка В.П.}


{О стационарных точках конформного 


радиуса спиралеобразных областей}{38}


\tl{Тронин С.Н., Семенова А.В.}


{Операды конечных помеченных графов}{50}


\tl{Уткина Е.А.} 


{О задачах Гурса с дополнительными нормальными производными 


в\newline краевых условиях}{61}


\tl{Филевич П.В.}


{К теореме Валирона о соотношениях


между максимумом модуля


и\newline максимальным членом целого ряда Дирихле}{66}


\tl{Эскин Л.Д.} 


{Уравнения, описывающие динамику неньютоновской жидкости с\newline реологическим 


законом Рейнера--Ривлина. \ II. \ Инвариантные решения}{73}














\vfill


\noindent\raisebox{2pt}{\copyright} Казанский госудаpственный унивеpситет


\end{document}


